\documentclass[11pt,a4paper]{article}

% ── Packages ──────────────────────────────────────────────────────────────────
\usepackage[margin=2.5cm]{geometry}
\usepackage{amsmath,amssymb}
\usepackage{graphicx}
\usepackage{booktabs}
\usepackage{longtable}
\usepackage{array}
\usepackage{xcolor}
\usepackage{hyperref}
\usepackage{caption}
\usepackage{subcaption}
\usepackage{enumitem}
\usepackage{microtype}
\usepackage{listings}
\usepackage{float}
\usepackage{placeins}

\graphicspath{{../figures/}}

\hypersetup{
  colorlinks = true,
  linkcolor  = blue!70!black,
  citecolor  = green!50!black,
  urlcolor   = blue!70!black,
}

% ── Custom commands ───────────────────────────────────────────────────────────
\newcommand{\numu}{$\nu_\mu$}
\newcommand{\numubar}{$\bar{\nu}_\mu$}
\newcommand{\uboone}{MicroBooNE}
\newcommand{\GeVc}{\,\text{GeV}\!/c}
\newcommand{\pot}{\,\text{POT}}
\newcommand{\cm}{\,\text{cm}}
\newcommand{\Npi}{$N\pi^\pm$}

\lstset{
  language=C++,
  basicstyle=\ttfamily\small,
  keywordstyle=\bfseries\color{blue!70!black},
  commentstyle=\itshape\color{gray},
  backgroundcolor=\color{gray!10},
  frame=single,
  framerule=0.4pt,
  numbers=none,
  breaklines=true,
}

\newcolumntype{C}[1]{>{\centering\arraybackslash}p{#1}}
\newcolumntype{R}[1]{>{\raggedleft\arraybackslash}p{#1}}
\newcolumntype{L}[1]{>{\raggedright\arraybackslash}p{#1}}
\emergencystretch=3em

% ─────────────────────────────────────────────────────────────────────────────
\begin{document}

\begin{titlepage}
  \centering
  \vspace*{2cm}
  {\LARGE\bfseries Measurement of $\nu_\mu+\bar\nu_\mu$ Charged-Current Multi-Pion Production
  on Argon with the MicroBooNE Detector in the NuMI Beam\par}
  \vspace{0.8cm}
  {\large Technical Note: the Two-Charged-Pion Final State
  (\texttt{CC1mu2piBDT})\par}
  \vspace{0.5cm}
  {\normalsize Draft 0.2, 2026-10-09\par}
  \vspace{1.5cm}
  % NOTE: author list to be confirmed. Contributors to the multi-pion effort:
  {\large Tusharadri Mahmud, Linhui Gu, Ishanee Pophale, Jarek Nowak\par}
  \vspace{0.5cm}
  {\normalsize\textit{Lancaster University}\par}
  \vspace{2cm}
  \begin{abstract}
    We describe the extension of the \uboone{} NuMI $\nu_\mu+\bar\nu_\mu$ charged-current
    single-charged-pion analysis to final states with two charged pions and any number of
    nucleons, CC$2\pi^\pm X p$. The selection follows the strategy of the BNB CC$2\pi^\pm Np$
    analysis~\cite{majeed}: a four-class particle classifier (muon, pion, proton, other) on every
    track, the assignment of the tracks to one muon and $N$ pions by the largest summed
    log-probability, and an event classifier on the assigned probabilities. On the full NuMI exposure
    ($7.766\times10^{20}$~POT in forward and $1.1082\times10^{21}$~POT in reverse horn current) it
    selects $13\%$ of the signal at $51\%$ purity in the combined configuration, against $18\%$ at
    $26\%$ for the selection of the first draft of this note. The uncertainty of the one-bin total is
    dominated by the flux, amplified by the subtraction of a background as large as the signal, and by
    the detector variations. Four control samples defined from the classifier outputs constrain the
    background through the conditional covariance of the extraction framework: in the combined
    configuration the one-bin total carries $15.7\%$ with the constraint and $46\%$ without, closes on
    pseudo-data and follows two reweighted interaction models. The detector-variation samples of the
    NuMI production were simulated with an earlier flux; the detector term is therefore taken
    fractionally. Three or more pions are not measurable at this exposure and serve as a control
    sample. The signal region is blind, and every result is extracted from pseudo-data.
  \end{abstract}
\end{titlepage}

\tableofcontents
\newpage

% =============================================================================
\section{Introduction}
\label{sec:intro}

Charged-current (CC) production of multiple charged pions on nuclei sits at the transition between
the resonant (RES) and deep-inelastic (DIS) regimes of neutrino--nucleus scattering and is strongly
shaped by the intranuclear hadronic cascade. A $\nu_\mu$ interaction that produces several pions at
the primary vertex may have some of them absorbed, charge-exchanged or produced in re-scattering
before they leave the nucleus, so the observed charged-pion multiplicity is a direct probe of
final-state interactions (FSI). Multi-pion final states are also a background to exclusive
single-pion and pion-less measurements, so a measurement of their rate on argon is useful both as a
cross-section result and as a constraint on the interaction model used across the \uboone{}
programme.

This note documents the extension of the NuMI $\nu_\mu+\bar\nu_\mu$ CC$1\pi^\pm$
analysis~\cite{note1} to two charged pions. The first draft of this note (2026-08-13) defined the
selections \texttt{CC1mu2pi} and \texttt{CC1mu3pi} as subselections of the single-pion selection
\texttt{CC1mu1piXp} with a looser pion identification~\cite{mahmud,gu}, on the Run-1 forward-horn
overlay only. Since then the analysis has been rebuilt on the strategy of the BNB CC$2\pi^\pm Np$
analysis~\cite{majeed} and on the current single-pion release: a baseline on the full exposure
(Section~\ref{sec:baseline}), a particle classifier (Section~\ref{sec:pid}), the track assignment and
an event classifier (Section~\ref{sec:event}), a study of the detector term
(Section~\ref{sec:detrobust}), control samples and a background constraint
(Section~\ref{sec:constraint}), and the extraction of the one-bin total through the framework
(Section~\ref{sec:framework}). The samples, normalisation, systematic universes and Wiener-SVD
unfolding are those of the single-pion analysis.

The decisions that fix the scope are:
\begin{itemize}[nosep]
  \item D1: the signal is inclusive in nucleons, CC$2\pi^\pm Xp$, with a subsample with at least one
        proton for comparison with the BNB analysis.
  \item D2: a per-pion momentum threshold of $0.10\GeVc$ for the signal, totals and angles;
        pion-momentum bins only above $0.175\GeVc$, where the range estimator is unbiased.
  \item D3: the two-pion signal region is blind. Part of it lies in the multi-$\pi$ control region of
        the single-pion analysis, which has been compared with beam-on data.
  \item D6: the truth requirement $\theta(\mu,\text{leading}~\pi)<2.6$~rad has a reconstructed
        counterpart on the assigned pions.
  \item D7: no exclusivity cut against the blind single-pion signal region; the overlapping events
        are opened only after the single-pion signal region.
  \item D8: the measurements stop at two pions. Events with three or more pion candidates are a
        control sample of the two-pion background.
\end{itemize}

% =============================================================================
\section{Signal definition}
\label{sec:signal}

The signal is a $\nu_\mu$ or $\bar\nu_\mu$ charged-current interaction with its vertex in the
fiducial volume of the single-pion analysis, one muon with $p_\mu>0.15\GeVc$ and exactly two charged
pions ($\pi^+$ or $\pi^-$; the pion charge is not measured), each with $p_\pi>0.10\GeVc$, no neutral
pion, kaon or heavier meson, and any number of nucleons. The opening angle between the muon and the
leading (hardest) pion is below $2.6$~rad, as in the single-pion definition. Pions are counted at any
momentum, so an event with a third charged pion below the threshold is not signal. The subsample with
at least one proton above $0.3\GeVc$ is $81\%$ of the signal.

\paragraph{Pion-momentum threshold.} In a two- or three-pion event the sub-leading pions are softer
than the leading one, and the softest pion sets how often the full topology can be reconstructed. On
the Run-1 forward-horn overlay the mean leading (softest) true pion momentum is $0.52$ ($0.27$)$\GeVc$
for two pions and $0.57$ ($0.21$)$\GeVc$ for three. The efficiency of the first-draft selection,
evaluated on the signal whose softest pion lies above a threshold $p_\pi^{\min}$, is given in
Table~\ref{tab:mp_thr}.

\begin{table}[H]
  \centering
  \caption{Signal efficiency of the first-draft selections as a function of a per-pion momentum
    threshold $p_\pi^{\min}$ applied to the softest pion, Run-1 forward-horn overlay. ``Kept'' is the
    fraction of the $p_\pi>0$ signal retained; Eff.\ is the efficiency within the restricted phase
    space.}
  \label{tab:mp_thr}
  \small
  \begin{tabular}{lcccc}
    \toprule
    & \multicolumn{2}{c}{Two pions} & \multicolumn{2}{c}{Three pions} \\
    \cmidrule(lr){2-3}\cmidrule(lr){4-5}
    $p_\pi^{\min}$ [$\GeVc$] & Kept & Eff.\ [\%] & Kept & Eff.\ [\%] \\
    \midrule
    $0$       & $1.00$ & $16.3$ & $1.00$ & $6.4$ \\
    $0.10$    & $0.90$ & $17.1$ & $0.82$ & $6.2$ \\
    $0.175$   & $0.70$ & $18.9$ & $0.49$ & $8.1$ \\
    $0.25$    & $0.43$ & $20.2$ & $0.26$ & $5.6$ \\
    \bottomrule
  \end{tabular}
\end{table}

Raising the threshold changes the efficiency little: with every pion above $0.175\GeVc$ about four in
five events are still lost, so the efficiency is limited by the reconstruction of several tracks from
one vertex, not by the soft pions. The threshold of $0.10\GeVc$, the pion tracking turn-on, removes the
unreconstructable tail ($10\%$ of the two-pion signal) and keeps $90\%$ of it; $0.175\GeVc$ would keep
$65\%$ of the two-pion signal in the combined configuration. Pion-momentum bins start at $0.175\GeVc$
(D2).

% =============================================================================
\section{Samples, exposure and normalisation}
\label{sec:samples}

The samples and normalisation are those of the single-pion release: fourteen $\nu_\mu$ overlay files
(forward horn current Runs 1, 2, 4, 5; reverse horn current Runs 1, 2, 3, 4a, 4b, 4c), seven beam-off
and seven dirt samples, each run normalised to its recorded exposure, $7.766\times10^{20}$~POT in
forward horn current (FHC), $1.1082\times10^{21}$~POT in reverse horn current (RHC) and
$1.8848\times10^{21}$~POT combined. The detector variations are the Run-4 FHC and RHC samples, a
central-value sample and eight variations each (light yield down and Rayleigh scattering,
recombination, space charge and four wire-response variations). The beam-on data are not used: the
pseudo-data are Poisson throws, run by run, of the overlays and the dirt at the data exposure, to
which the framework adds the beam-off prediction. Generator predictions with the two-pion signal
definition have not yet been produced.

% =============================================================================
\section{Baseline on the full exposure}
\label{sec:baseline}

\texttt{CC1mu2pi} and \texttt{CC1mu3pi} inherit the \texttt{CC1mu1piXp} selection (slice, fiducial
volume, topology and muon candidate) with a looser pion identification: vertex distance up to
$9.5\cm$, the pion BDT of the first draft in place of the single-pion identification, up to two
uncontained pion candidates, no shower veto and no single-pion opening-angle cut. Following D6 they
now require $\theta(\mu,\text{longest pion candidate})<2.6$~rad. The selections store every counted
candidate (containment, length, range and MCS momenta, direction, identification scores and the
backtracked truth), the true pions, the proton and kaon counts and a truth topology class. On a
30\,000-event slice every branch of the single-pion selection is identical before and after these
changes. Table~\ref{tab:baseline} gives the baseline at the full exposure, beam-off and dirt included.

\begin{table}[H]
  \centering
  \caption{Baseline of the first-draft selections at the full exposure, with the reconstructed
    opening-angle requirement. Purities include beam-off and dirt.}
  \label{tab:baseline}
  \small
  \begin{tabular}{lrrrcc}
    \toprule
    & Generated signal & Selected signal & Selected total & Efficiency & Purity \\
    \midrule
    Two pions, FHC       & $1\,210$ & $221$ & $836$     & $18.2\%$ & $26.4\%$ \\
    Two pions, RHC       & $1\,672$ & $303$ & $1\,168$  & $18.1\%$ & $26.0\%$ \\
    Two pions, combined  & $2\,882$ & $524$ & $2\,004$  & $18.2\%$ & $26.1\%$ \\
    \midrule
    Three pions, FHC      & $204$ & $21.7$ & $203$ & $10.6\%$ & $10.7\%$ \\
    Three pions, RHC      & $347$ & $32.8$ & $314$ & $9.5\%$  & $10.4\%$ \\
    Three pions, combined & $551$ & $54.5$ & $517$ & $9.9\%$  & $10.5\%$ \\
    \bottomrule
  \end{tabular}
\end{table}

These numbers replace the raw Run-1 counts of the first draft, whose ``670 (66) events at full
exposure'' were overlay counts at the simulated exposure. The largest two-pion backgrounds are
single-pion events with a second counted candidate ($22\%$ of the sample), events with a neutral pion
($21\%$), three or more charged pions ($7\%$), neutral-current events ($6\%$) and beam-off ($2.5\%$).
In selected signal events only $60\%$ have both counted candidates backtracked to true pions, and the
longest candidate is the leading pion in $52\%$, so observables ordered by pion cannot be used with
this candidate choice. The three-pion sample is $10.5\%$ pure and its largest background, $30\%$, is a
two-pion event with a third counted candidate.

% =============================================================================
\section{Particle classifier}
\label{sec:pid}

A four-class classifier (muon, pion, proton, other) is trained with XGBoost on the tracks of the
pion-candidate pool of the selection (generation 2, track score $\geq0.3$), dumped from all fourteen
overlays ($3.1$~million tracks). Tracks are labelled by their backtracked particle; tracks without a
truth match, mostly cosmic rays of the overlay, are the class ``other''. Training uses run periods 1,
2 and 4 of both horn modes; the held-out periods are RHC Run~3 and FHC Run~5. The beam-on, beam-off and
dirt ntuples of the older production (Runs 1--3) lack twenty of the track branches every overlay
carries (the Bragg-pion likelihoods, truncated $dE/dx$, deflection, end spacepoints, hit and
descendant counts), so the deployed models use the 43 inputs every sample carries. The model is
evaluated in the selection through ROOT's \texttt{RBDT}, which reproduces XGBoost exactly.

\begin{table}[H]
  \centering
  \caption{ROC area of pion identifications on the two-pion candidate pool, held-out runs, tracks with
    a truth match (backtracked purity above $0.5$).}
  \label{tab:pidroc}
  \small
  \begin{tabular}{lccc}
    \toprule
    Pion identification & $\pi$ vs rest & $\pi$ vs $p$ & $\pi$ vs $\mu$ \\
    \midrule
    Particle classifier (63 inputs)          & $0.954$ & $0.985$ & $0.895$ \\
    Pion BDT of the first draft, as applied   & $0.888$ & $0.945$ & $0.804$ \\
    Pion BDT of the single-pion analysis     & $0.818$ & $0.961$ & $0.558$ \\
    MIP BDT                                  & $0.718$ & $0.957$ & $0.271$ \\
    LLR score                                & $0.670$ & $0.893$ & $0.251$ \\
    \bottomrule
  \end{tabular}
\end{table}

The gain is the same in FHC and RHC and for $\pi^+$ and $\pi^-$. With the unmatched tracks as the
fourth class, at the pion efficiency of the first-draft BDT ($91.8\%$) the purity of the counted pion
candidates rises from $27.6\%$ to $43.3\%$ and the fraction of protons passing falls from $13.4\%$ to
$4.1\%$ (63 inputs). With the 43 inputs every sample carries, the pion ROC area on the full pool,
unmatched tracks included, is $0.938$ ($0.946$ with all inputs, $0.868$ for the first-draft BDT). The
pion efficiency does not depend on the true pion momentum more than that of the first-draft BDT and is
higher at threshold. Muon against pion is the hardest separation ($0.895$; about a third of the true
muons in the pool still pass).

Used as the pion identification of the single-pion selection at the same pion efficiency, the
43-input classifier raises the purity of the counted tracks from $52.6\%$ to $64.8\%$ and, at event level on the
held-out runs, the selection purity from $58.0\%$ to $64.7\%$ for $5\%$ of the signal. This is a study
for a later revision of the single-pion analysis.

\paragraph{The LLR-and-geometry classifier.} Section~\ref{sec:detrobust} shows that the detector
variations act on the calorimetric inputs. A second model drops the Bragg likelihoods, $\chi^2$
scores, PIDA and calorimetric energies, and keeps the LLR score (summed and per plane), length, track
score, vertex distance, start and end containment, range and MCS momenta, the proton-hypothesis energy,
the daughter counts and the per-plane hit flags. Its pion ROC area is $0.892$ against $0.926$ for the
deployed model on the same tracks, and at $80\%$ pion efficiency it passes $9.0\%$ of the protons
against $3.9\%$. It is the particle classifier of \texttt{CC1mu2piBDT}.

% =============================================================================
\section{Track assignment and event classifier}
\label{sec:event}

\paragraph{Assignment.} In events with one neutrino slice and the vertex in the fiducial volume, the
tracks of the pion pool with track score $\geq0.5$, the muon candidate included, carry the four class
probabilities. For the hypothesis ``one muon, $N$ pions, the other tracks protons or other'' the
log-likelihood is
\begin{equation}
  L_N = \log P_\mu(\mu) + \sum_{\text{pions}}\log P_\pi + \sum_{\text{rest}}\log\left(1-P_\mu-P_\pi\right),
\end{equation}
and every choice of the muon and the $N$ pions among at most eight tracks is enumerated; the best and
the runner-up are kept for $N=1,2,3$. In the selected signal the assigned muon and pions are correct
in $91$--$92\%$ of events. In the preselected signal, however, only $30\%$ of the events have the muon
and both pions as tracks, so in most signal events no assignment can be fully right.

\paragraph{Event classifier.} An XGBoost classifier separates the two-pion signal from everything else
on 33 features: the probabilities of the assigned tracks, $L_2-L_1$, $L_2-L_3$, the margin over the
runner-up, the class sums and the counts above $0.5$ over all tracks, the track and shower counts, the
topological score and CosmicIP. Its domain is one slice, the vertex in the fiducial volume, the software
trigger, at least three tracks and $\theta(\mu,\text{longest assigned pion})<2.6$~rad (D6). It is trained
on run periods 1, 2 and 4 of both horn modes ($206\,152$ events, $13\,558$ signal) with the nominal
mixture of overlay, dirt and beam-off and the central-value weights, and tested on FHC Run~5 and RHC
Run~3 ($104\,814$ events, $7\,253$ signal), which the particle classifier did not see either. The ROC
area on the test periods is $0.887$; adding the lengths, vertex distances, containment and opening
angles of the assigned tracks raises it to $0.898$, a gain that disappears in the total uncertainty, so
these kinematic features are not used. The largest gains come from the summed pion probability, $L_2$,
the number of tracks with $P_\pi>0.5$, the topological score and the number with $P_\mu>0.5$.

\paragraph{Working point.} The threshold is chosen on the expected total uncertainty of the one-bin
total, with the test-period yields projected to the full exposure: data statistics, flux
($17\%\times(1+B/S)$), POT and targets, cross section, re-interaction, detector (the eight Run-4
variations, FHC and RHC correlated) and simulation statistics (Table~\ref{tab:wp}). With statistics and
flux alone the minimum lies at $3.9\%$ efficiency; the detector term rises steeply above a score of
$0.95$, and with all terms the minimum is at $0.92$, flat between $0.90$ and $0.94$ within the
statistics of the detector-variation samples.

\begin{table}[H]
  \centering
  \caption{Expected uncertainty of the one-bin total (combined configuration) for the first-draft
    selection, the same with the particle classifier as pion identification
    (\texttt{CC1mu2piNewPID}), and the event classifier with the deployed and with the
    LLR-and-geometry particle classifier at score $>0.92$. The last row is the evaluation of
    Section~\ref{sec:constraint}, with the flux from the PPFX universes instead of
    $17\%\times(1+B/S)$.}
  \label{tab:wp}
  \footnotesize\setlength{\tabcolsep}{4pt}
  \begin{tabular}{lcccccccc}
    \toprule
    Selection & Eff. & Purity & Stat. & Flux & Cross sect. & Reint. & Detector & Total \\
    \midrule
    \texttt{CC1mu2pi}                     & $18.2\%$ & $26\%$ & $8\%$ & $65\%$ & $47\%$ & $17\%$ & $34\%$ & $89\%$ \\
    \texttt{CC1mu2piNewPID}               & $15.8\%$ & $35\%$ & $8\%$ & $48\%$ & $27\%$ & $13\%$ & $21\%$ & $61\%$ \\
    Event classifier, deployed            & $17.1\%$ & $56\%$ & $6\%$ & $31\%$ & $12\%$ & $5\%$  & $33\%$ & $47\%$ \\
    Event classifier, LLR and geometry    & $13.0\%$ & $51\%$ & $7\%$ & $35\%$ & $17\%$ & $6\%$  & $21\%$ & $46\%$ \\
    \bottomrule
  \end{tabular}
\end{table}

The cross-section term is the model dependence of the background and falls with the purity; the model
dependence of the efficiency stays below $6\%$. In the combined configuration the signal region of the
LLR-and-geometry selection holds $738$ events, $378$ of them signal; the largest backgrounds are
single-pion events ($15\%$ of the sample) and three or more charged pions ($13\%$). Of the selected
events $14\%$ lie in the blind single-pion signal region (D7) and $12\%$ in the opened multi-$\pi$
control region of the single-pion analysis (D3).

\paragraph{Selection requirements.} Table~\ref{tab:cuts} lists the requirements of
\texttt{CC1mu2piBDT}. The cuts of the single-pion selection after its preselection (topological score
above $0.67$ and the muon-candidate, pion-containment, wire-gap, shower and track-multiplicity
requirements) are not applied; the topological score, CosmicIP and the track and shower counts enter
only as features of the event classifier.

\begin{table}[H]
  \centering
  \caption{Requirements of \texttt{CC1mu2piBDT}. Cuts 1--9 define the domain of the event classifier;
    the signal region adds cut 10. The sidebands (Section~\ref{sec:constraint}) are the events that
    pass cuts 1--9 and fail cut 10, assigned in the order listed; in the extraction they exclude the
    events selected by the single-pion analysis.}
  \label{tab:cuts}
  \small
  \begin{tabular}{rlL{9.4cm}}
    \toprule
    & Cut & Requirement \\
    \midrule
    \multicolumn{3}{l}{\emph{Event}} \\
    1  & Neutrino slices     & exactly one \\
    2  & Reconstructed vertex & inside the fiducial volume of the single-pion analysis, outside the
                               dead-wire region in $z$ \\
    3  & Software trigger    & passed \\
    \multicolumn{3}{l}{\emph{Track pool: direct neutrino daughters}} \\
    4  & Track score         & $\geq0.5$ \\
    5  & PID inputs          & LLR score in $(-1,2)$; proton, muon and MIP Bragg likelihoods in
                               $(0,500)$; length $>0$ \\
    6  & Ranking             & by $\max(P_\mu,P_\pi)$ of the LLR-and-geometry particle classifier; at
                               most eight tracks kept \\
    7  & Pool size           & at least three tracks \\
    \multicolumn{3}{l}{\emph{Assignment and event classifier}} \\
    8  & Assignment          & the muon and two pions that maximise $L_2$ \\
    9  & Opening angle       & $\theta(\mu,\text{longer assigned pion})<2.6$~rad (D6) \\
    10 & Event score         & $>0.92$: signal region \\
    \midrule
    \multicolumn{3}{l}{\emph{Sidebands: cuts 1--9, score $\leq0.92$}} \\
       & Three-pion-like     & at least four tracks and $L_3>L_2$ \\
       & Proton-like         & an assigned pion with $P_p>0.5$ \\
       & Score band          & $0.5<\text{score}\leq0.92$ \\
       & Low score           & score $\leq0.5$ \\
    \bottomrule
  \end{tabular}
\end{table}

% =============================================================================
\section{The detector term}
\label{sec:detrobust}

At the working point of the deployed classifier the detector term, $33\%$, is the largest. A paired
bootstrap over the physical events, with one Poisson weight per event shared by the central-value
sample and the eight variations, attributes $7.7\%$ of it to the statistics of the variation samples.
It comes from the recombination ($+22\%$) and wire-response variations (angle $yz$ $+18\%$, $x$
$+12\%$, $yz$ $+8\%$); the light-yield, space-charge and angle-$xz$ variations are consistent with
zero. At track level these variations move the $\chi^2$ scores, PIDA, the Bragg likelihoods and the
calorimetric energies by up to half of their spread for protons, the LLR score by at most $0.04$ and the
geometric inputs not at all. At event level they lower the selected signal by about $8\%$ and the
background by $10$--$18\%$, and the detector term splits into an efficiency part ($14.3\%$), which no
background constraint can remove, and a background part ($17.8\%$).

\begin{table}[H]
  \centering
  \caption{Variants of the particle classifier and of the event features (combined configuration), at
    the best threshold of each. The detector term is split into its efficiency and background parts.}
  \label{tab:detvariants}
  \footnotesize\setlength{\tabcolsep}{4pt}
  \begin{tabular}{llcccccc}
    \toprule
    Particle classifier & Event features & Score $>$ & Eff. & Purity & Detector & (eff., bkg.) & Total \\
    \midrule
    Deployed                       & PID set  & $0.92$ & $17.1\%$ & $56\%$ & $32.7\%$ & ($14.3$, $17.8$) & $47\%$ \\
    Deployed                       & full set & $0.94$ & $14.4\%$ & $62\%$ & $33.0\%$ & ($18.2$, $14.8$) & $45\%$ \\
    Without $\chi^2$, PIDA, energy & PID set  & $0.93$ & $14.4\%$ & $57\%$ & $35.1\%$ & ($14.9$, $19.0$) & $48\%$ \\
    LLR and geometry               & PID set  & $0.92$ & $13.0\%$ & $51\%$ & $21.4\%$ & ($7.8$, $15.2$)  & $44\%$ \\
    LLR and geometry               & full set & $0.93$ & $12.0\%$ & $55\%$ & $20.3\%$ & ($8.4$, $12.9$)  & $42\%$ \\
    \bottomrule
  \end{tabular}
\end{table}

No variant removes the detector term, and the totals lie within the $3\%$ scatter of the scans. The
LLR-and-geometry classifier halves the efficiency part ($14.3\%$ to $7.8\%$), the part a constraint
cannot reach, and is adopted. Training the event classifier or the particle classifier with the
variation samples did not help.

For comparison, the BNB analysis quotes a detector term of $11.0\%$ for the predicted selected event
count, from seven Run-4d variations without the wire-angle ones~\cite{majeed}. In the same convention
the LLR-and-geometry selection has $10.6\%$ with all eight variations and $8.4\%$ with those seven; the
terms of Table~\ref{tab:wp} are those of the extracted total, in which the background subtraction and
the efficiency correction amplify the change of the predicted events by a factor of two to four. The
recombination variation leads in both analyses.

% =============================================================================
\section{Control samples and the background constraint}
\label{sec:constraint}

The flux term scales the signal and the background together, and the background is as large as the
signal. Control samples that share the background composition constrain it through the conditional
covariance of the extraction framework (\texttt{ConstrainedCalculator}): the covariance of the
predicted counts in the signal region and the sidebands, built from every source, is conditioned on the
sidebands, and the background estimate moves with the sideband data. Four sidebands are defined from
the classifier outputs, exclusive and in this order inside the two-pion domain
(Table~\ref{tab:cuts}): three-pion-like (at
least four tracks and three pions more likely than two; the control sample of D8), proton-like (an
assigned pion with $P_p>0.5$), event score between $0.5$ and the working point, and event score below
$0.5$. The sidebands leave out the events of the blind single-pion signal region, so they can be opened
without it (Table~\ref{tab:regions}).

\begin{table}[H]
  \centering
  \caption{Signal region and sidebands of the LLR-and-geometry selection, combined configuration, full
    exposure.}
  \label{tab:regions}
  \small
  \begin{tabular}{lrrcl}
    \toprule
    Region & Events & Signal & Purity & Largest classes \\
    \midrule
    Signal region (score $>0.92$)   & $738$       & $378$ & $51\%$ & CC$1\pi$ $15\%$, CC$\geq3\pi$ $13\%$ \\
    Three-pion-like                 & $282$       & $39$  & $14\%$ & CC$\geq3\pi$ $26\%$, CC$\pi^0$ $18\%$, beam-off $15\%$ \\
    Proton-like                     & $15\,160$   & $431$ & $3\%$  & CC$0\pi$ $40\%$, CC$\pi^0$ $16\%$, CC$1\pi$ $15\%$ \\
    Score $0.5$--$0.92$             & $4\,397$    & $669$ & $15\%$ & CC$\pi^0$ $27\%$, CC$1\pi$ $21\%$ \\
    Score $<0.5$                    & $10\,166$   & $143$ & $1\%$  & beam-off $45\%$, CC$0\pi$ $11\%$, CC$\pi^0$ $11\%$ \\
    \bottomrule
  \end{tabular}
\end{table}

A first evaluation outside the framework, with the formalism of \texttt{ConstrainedCalculator},
gives the uncertainties of Table~\ref{tab:constraint}. The flux term falls from $35\%$ to $5.5\%$, since
the flux scales the background and the sidebands alike, and the detector term from $21\%$ to $9\%$,
close to its efficiency part. The three-pion-like sample alone, $282$ events, halves the total.

\begin{table}[H]
  \centering
  \caption{Uncertainty of the one-bin total (combined configuration) with each sideband constraint,
    evaluated outside the framework (flux from the PPFX universes, detector taken fractionally).}
  \label{tab:constraint}
  \small
  \begin{tabular}{lccccccc}
    \toprule
    Constraint & Total & Flux & Cross sect. & Reint. & Detector & Data stat. & MC stat. \\
    \midrule
    None              & $45.5\%$ & $35.4\%$ & $16.7\%$ & $5.5\%$ & $20.6\%$ & $7.2\%$ & $3.9\%$ \\
    Three-pion-like   & $24.2\%$ & $13.6\%$ & $8.3\%$  & $4.2\%$ & $11.7\%$ & $11.4\%$ & $6.8\%$ \\
    Proton-like       & $31.9\%$ & $13.0\%$ & $15.2\%$ & $5.3\%$ & $22.9\%$ & $7.3\%$ & $4.0\%$ \\
    Score $0.5$--$0.92$ & $22.5\%$ & $6.2\%$ & $8.4\%$  & $5.2\%$ & $17.2\%$ & $7.7\%$ & $4.2\%$ \\
    Score $<0.5$      & $33.0\%$ & $13.6\%$ & $10.6\%$ & $4.9\%$ & $26.2\%$ & $7.7\%$ & $4.4\%$ \\
    All sidebands     & $16.8\%$ & $5.5\%$  & $7.0\%$  & $4.8\%$ & $8.8\%$  & $8.8\%$ & $5.0\%$ \\
    \bottomrule
  \end{tabular}
\end{table}

\paragraph{Closure under alternative models.} Pseudo-data from a reweighted interaction model in every
region, extracted with the nominal model, return the alternative truth: with all sidebands $0.952$
($-0.29\sigma$) under GENIE multisim universe~545 and $0.994$ ($-0.03\sigma$) under the
$\Delta\to N\pi$ angular variation, against $0.979$ and $1.012$ without the constraint; over the GENIE
multisim universes the bias has an rms of $4.5\%$ with the constraint, within $1\sigma$ in every universe,
and $17.4\%$ without ($98\%$ of the universes within $1\sigma$). These models lie inside the variations that build the covariance, so the test checks
the procedure, not the model space. The score band holds more signal than the signal region ($669$
against $378$ events), so its constraint carries signal-model information into the background
estimate; it closes for both models.

% =============================================================================
\section{Extraction through the framework}
\label{sec:framework}

\texttt{CC1mu2piBDT} implements the selection in the framework: the track pool, the class
probabilities of the LLR-and-geometry classifier, the assignment, the 33 features, the event score
(\texttt{RBDT}) and the region; the selected events are those of the signal region. On a
30\,000-event slice the features equal those of the Python study exactly, the scores to $6\times10^{-8}$
and the regions in every event, and \texttt{CC1mu1piXp} is bitwise unchanged. The 28 release samples
and the 18 detector-variation samples were processed with it. The bin configuration has the signal
region as an ordinary bin and the four sidebands as sideband bins in the signal block, so that their
signal is predicted. The keyword \texttt{Calculator Constrained} in the extraction configuration selects
\texttt{ConstrainedCalculator}; the extractor propagates the ordinary block of each covariance and,
with the constraint, the part that remains after the conditioning, so that the terms by source add up
to the constrained total.

\begin{table}[H]
  \centering
  \caption{One-bin two-pion total through the framework, on pseudo-data, without and with the sideband
    constraint; detector variations taken fractionally and weighted to the current flux
    (Section~\ref{sec:detvar}). The last column is the extracted value over the truth of the pseudo-data.}
  \label{tab:framework}
  \footnotesize\setlength{\tabcolsep}{4pt}
  \begin{tabular}{lccccccc}
    \toprule
    & Total & Detector & Flux & Cross sect. & Reint. & Data stat. & Unfolded/truth \\
    \midrule
    FHC, no constraint       & $54.0\%$ & $34.2\%$ & $37.0\%$ & $14.7\%$ & $5.2\%$ & $10.5\%$ & $1.056$ \\
    FHC, constrained         & $22.5\%$ & $11.4\%$ & $6.1\%$  & $8.7\%$  & $5.9\%$ & $14.0\%$ & $1.020$ \\
    RHC, no constraint       & $41.1\%$ & $18.1\%$ & $31.5\%$ & $15.5\%$ & $5.4\%$ & $8.7\%$  & $1.035$ \\
    RHC, constrained         & $16.3\%$ & $7.3\%$  & $4.9\%$  & $5.9\%$  & $4.1\%$ & $11.1\%$ & $1.020$ \\
    Combined, no constraint  & $46.0\%$ & $25.9\%$ & $33.7\%$ & $15.0\%$ & $5.2\%$ & $6.7\%$  & $1.044$ \\
    Combined, constrained    & $15.7\%$ & $8.2\%$  & $4.8\%$  & $6.6\%$  & $5.0\%$ & $8.8\%$  & $1.021$ \\
    \bottomrule
  \end{tabular}
\end{table}

The constrained combined total, $15.7\%$, agrees with the evaluation outside the framework
($16.8\%$). With the constraint the data statistics become the largest term, and the three
configurations close at $1.020$--$1.021$.

\subsection{Detector-variation samples}
\label{sec:detvar}

The first framework extraction gave a larger detector term than the study outside it, most in RHC
($28\%$ against $17\%$). Per POT, the Run-4 detector-variation central-value samples predict $1.27$
(FHC) and $1.51$ (RHC) times the event rate of the nominal Run-4 overlays in the fiducial volume, the
ratio rising from $1.1$ below $0.5$~GeV to $1.8$--$3$ above $5$~GeV. The samples were simulated with the
earlier NuMI flux, the nominal overlays with the current one (Geant~v4.10.4): weighted event by event
with the ratio of the two fluxes in true energy and angle to the beamline, they predict $0.94$ times
the nominal rate, flat in energy. The framework built the detector covariance from the absolute
difference between each variation and the variation central value, and so carried that ratio. The
covariance type \texttt{DVfrac} applies the fractional change of each variation relative to the
variation central value to the nominal prediction; it is used here and, since version~1.8, in the
single-pion release.

Because the earlier flux overweights high energies, where the two-pion signal sits, the variation
samples are in addition weighted to the current flux: the selection framework applies an optional
per-event weight branch, carried only by reweighted copies of the 18 variation samples. On the
fractional detector change of the selected two-pion count the weight raises the FHC term from $17.6\%$
to $20.6\%$ and leaves RHC unchanged ($9.3\%$ to $9.1\%$); for the single-pion selection it changes the
term by at most $0.3\%$ absolute. Table~\ref{tab:detvar} shows the three treatments.

\begin{table}[H]
  \centering
  \caption{Detector term and total of the one-bin two-pion total for the three treatments of the
    detector variations.}
  \label{tab:detvar}
  \small
  \begin{tabular}{lcccccc}
    \toprule
    & \multicolumn{2}{c}{Absolute} & \multicolumn{2}{c}{Fractional} & \multicolumn{2}{c}{Fractional, flux-weighted} \\
    \cmidrule(lr){2-3}\cmidrule(lr){4-5}\cmidrule(lr){6-7}
    & Detector & Total & Detector & Total & Detector & Total \\
    \midrule
    FHC, no constraint       & $40.8\%$ & $58.4\%$ & $29.8\%$ & $51.3\%$ & $34.2\%$ & $54.0\%$ \\
    FHC, constrained         & $15.4\%$ & $24.9\%$ & $11.9\%$ & $22.4\%$ & $11.4\%$ & $22.5\%$ \\
    RHC, no constraint       & $28.4\%$ & $46.6\%$ & $17.9\%$ & $41.0\%$ & $18.1\%$ & $41.1\%$ \\
    RHC, constrained         & $11.1\%$ & $18.5\%$ & $7.6\%$  & $16.3\%$ & $7.3\%$  & $16.3\%$ \\
    Combined, no constraint  & $33.4\%$ & $50.7\%$ & $23.0\%$ & $44.5\%$ & $25.9\%$ & $46.0\%$ \\
    Combined, constrained    & $11.9\%$ & $17.9\%$ & $8.5\%$  & $15.7\%$ & $8.2\%$  & $15.7\%$ \\
    \bottomrule
  \end{tabular}
\end{table}

The sideband constraint removes the difference between the fractional treatments: the constrained totals
are the same with and without the flux weight.

% =============================================================================
\section{Three or more pions}
\label{sec:threepi}

In the two-pion domain, events with at least four tracks in which three pions are more likely than two
number $234$ at full exposure (combined, deployed particle classifier), of which $88$ have three or more charged pions, $41$ a neutral
pion and $35$ are two-pion signal. A three-pion event classifier built the same way reaches a ROC area of
$0.855$, but with all terms every three-pion selection carries a total uncertainty above $100\%$
(\texttt{CC1mu3pi} $245\%$, with the particle classifier $117\%$, the event classifier at best
$105\%$). Four and five pions are out of reach: combined, $3.1$ selected four-pion signal events at
$4.8\%$ purity and $0.1$ five-pion events, and the probability that a true pion multiplicity is
reconstructed in its own bin is $0.84$, $0.48$ and $0.33$ for one, two and three or more pions. The
three-pion-like events are therefore the control sample of Section~\ref{sec:constraint} (D8).

% =============================================================================
\section{Status and next steps}
\label{sec:status}

\begin{itemize}[nosep]
  \item The two-pion one-bin total is extracted through the framework in FHC, RHC and combined, with
        and without the sideband constraint (Table~\ref{tab:framework}).
  \item The LLR score is mis-modelled in the opened control regions of the single-pion analysis (shape
        $\chi^2$ of $160$--$250$ for 19 degrees of freedom, mean shifts of $-0.007$ to $-0.031$), and
        the particle classifier relies on it. The inputs of the LLR-and-geometry classifier and the event
        score are to be compared with beam-on data in the opened control regions, and the effect on the
        sidebands quantified, before the sidebands are opened.
  \item Validation through the framework: ensembles of pseudo-experiments, the alternative-model
        closure with the constraint, and an independent-generator closure if a NuMI NuWro overlay can
        be validated.
  \item Generator predictions with the two-pion signal definition.
  \item Response and binning of the differential candidates ($p_\mu$, $\cos\theta_\mu$, pion angles
        with a convention symmetric in the pions, the opening angle between the pions, the proton
        multiplicity) under the $0.50$ migration criterion; the invariant masses remain study products
        unless they pass it.
  \item The same sideband constraint for the single-pion analysis, on pseudo-data.
\end{itemize}

% =============================================================================
\begin{thebibliography}{9}
\bibitem{mahmud} T.\ Mahmud, \textit{Study of $\nu_\mu$ charged-current three-charged-pion
  production in \uboone{}} (\texttt{CC3}$\pi^+$ selection), MSc by Research thesis,
  Lancaster University.
\bibitem{gu} L.\ Gu, multi-pion event-selection analyzer for the \uboone{} NuMI
  cross-section framework (private communication / internal code).
\bibitem{majeed} N.\ Majeed, \textit{Study of production mechanisms of CC events with two charged pions
  and at least one proton in the final state}, \uboone{} internal note, version~2 (May 2026).
\bibitem{note1} I.\ Pophale and J.\ Nowak, \textit{Measurement of
  $\nu_\mu+\bar\nu_\mu$ charged-current single-pion production differential cross sections
  on argon with \uboone{} in the NuMI beam}, \uboone{} analysis note, version~1.8 (companion document).
\end{thebibliography}

\end{document}
